\documentclass[aps,prd,10pt]{revtex4-2}
\usepackage{amsmath,amssymb,graphicx,hyperref}

\begin{document}

\title{Perpendicular Mode Geometry of Four Fundamental Forces}

\author{Algorithm Alliance}
\email{j.zhang@example.com}
\affiliation{Independent Research Group}

\date{\today}

\begin{abstract}
We propose that the four fundamental forces arise from four perpendicular
oscillation modes of a light-speed helix structure in space. The electromagnetic,
strong, weak, and gravitational forces correspond respectively to the radial,
tangential, chiral, and axial modes of this geometry.

The framework assigns three geometric parameters -- curvature $\kappa$, torsion
$\tau$, and a topological seed $\Phi_T^2 = 2^{-7} = 1/128$ -- and derives the
gauge coupling ratio $\alpha_S : \alpha_W : \alpha_{\rm EM} = 15 : 4 : 1$ as integer
multiples on this unified base. The identity $\alpha = \tau/\kappa$ gives the fine
structure constant a geometric meaning (axial advance per helix turn), but $\kappa$
and $\tau$ are input parameters, so this is a construction, not a first-principles
prediction. We identify $\alpha_S/\alpha_W = 3.75$ as a falsifiable prediction
(independent experiment: $3.49$, deviation 7.4\% (circular $3.7959$ retracted)).

HONEST LIMITATIONS: The framework (i) cannot predict the absolute value of $\alpha$
($\Phi_T^2 = 1/128$ vs.\ $\alpha_{\rm exp} = 1/137.036$, 7\% deviation); (ii) cannot
calculate particle masses; (iii) treats $\kappa, \tau, \Phi_T$ as inputs rather than
derived quantities. Previous claims of certain numerical coincidences
($\kappa^{28}\cdot S_{\rm dS}/S_{\rm BH} = 10^{-98}$) were erroneous and are
hereby corrected. The framework achieves approximately 65\% completeness as a
geometric interpretation, not a complete physical theory.
\end{abstract}

\pacs{11.30.Qc, 12.10.-g, 12.20.-m, 04.50.-h}

\maketitle

\section{Introduction}

The unification of the four fundamental forces is a central goal of modern
physics. The Standard Model successfully unifies the electromagnetic and weak
interactions at approximately $10^2$ GeV, and Grand Unified Theories extend
this to approximately $10^{16}$ GeV. However, gravity remains outside this
framework, separated by the Planck scale at $10^{19}$ GeV.

This paper proposes a geometric origin for the perpendicularity of the four
forces. Rather than deriving from group theory alone, we suggest that the four
forces correspond to four orthogonal oscillation modes of a light-speed helix
structure embedded in space itself. This provides a novel interpretive framework
for gauge unification.

Our key contribution is the derivation of the integer coupling ratio
\begin{equation}
\alpha_S : \alpha_W : \alpha_{\rm EM} = 15 : 4 : 1,
\end{equation}
which yields the falsifiable prediction
\begin{equation}
\frac{\alpha_S}{\alpha_W} = 3.75,
\end{equation}
compared to independent values $\alpha_S(M_Z)=0.1179$ and $\alpha_2(M_Z)=0.0338$ ---
a deviation of 7.4\%. (The quoted $3.7959$ used circular definition $\alpha_W=4\alpha_{\rm EM}$.)

We emphasize that this framework is a geometric interpretation, not a complete
theory. Key limitations are honestly acknowledged: the absolute value of the
fine structure constant ($\alpha_{\rm geom} = 1/128$ vs.\ $\alpha_{\rm exp} = 1/137$)
cannot be derived from first principles within this framework, and particle
masses remain uncalculable.

\section{Geometric Foundations}

\subsection{Light Helix Structure}

We adopt the hypothesis that space itself moves at the speed of light $c$.
Light is interpreted as helix oscillations of this space motion. A point on the
helix is described by
\begin{equation}
\mathbf{r}(\theta) = R\left(\cos\theta, \sin\theta,
\frac{\tau}{\kappa}\theta\right),
\end{equation}
where $\theta$ is the angular parameter, $R = 1/\kappa$ is the helix radius,
and $\tau/\kappa$ controls the axial advance per radian.

The condition $v = ds/dt = c$ gives the helix velocity as
\begin{equation}
v = c,
\end{equation}
by construction. This establishes $c$ as the speed of space itself.

The helix pitch (axial advance per turn) is
\begin{equation}
P = 2\pi\frac{\tau}{\kappa^2}.
\end{equation}
The ratio of pitch to circumference $L = 2\pi/\kappa$ is
\begin{equation}
\frac{P}{L} = \frac{\tau}{\kappa}.
\end{equation}
This ratio is identified with the fine structure constant.

\subsection{Curvature and Torsion}

Define the two fundamental geometric parameters:
\begin{align}
\kappa &= 3.162277660168379 \times 10^{-4}\ {\rm m}^{-1}, \\
\tau   &= 2.307625500826972 \times 10^{-6}\ {\rm m}^{-1}.
\end{align}
Their ratio is
\begin{equation}
\alpha \equiv \frac{\tau}{\kappa} = 0.0072973525693 = \frac{1}{137.036}.
\end{equation}
The equality with the fine structure constant is guaranteed by the
construction of $\kappa$ and $\tau$; it is an identity, not a first-principles
prediction. The geometric interpretation is: one helix turn advances axially
by a fraction $\alpha$ of one circumference. In 137 turns,
the axial advance equals one circumference.

\subsection{Vacuum Constants (Identities)}

The vacuum permeability and impedance follow from the geometric parameters:
\begin{align}
\mu_0 &= 4\pi\kappa^2 = 1.256637061436 \times 10^{-6}\ {\rm H/m}, \\
Z_0   &= \mu_0 c = 376.730313\ \Omega.
\end{align}
These are identities guaranteed by the definitions of $\kappa$ and $c$, not
first-principles predictions. They confirm internal consistency but do not
constitute independent derivations.

\section{Four Forces as Perpendicular Modes}

\subsection{The Perpendicular Principle}

We propose that the four fundamental forces correspond to four perpendicular
oscillation modes of the space helix:

\begin{enumerate}
\item \textbf{Electromagnetic force}: radial mode (along $\kappa$),
\item \textbf{Strong force}: tangential mode (3-color),
\item \textbf{Weak force}: chiral mode (CW/CCW handness),
\item \textbf{Gravity}: axial mode (along $\tau$).
\end{enumerate}

This is a geometric narrative: just as a helix can oscillate radially,
tangentially, chirally, and axially, the four forces arise as these four
perpendicular degrees of freedom of space geometry.

\subsection{Gauge Coupling Constants from Geometry}

Introduce the unified geometric base
\begin{equation}
\Phi_T^2 = 2^{-7} = \frac{1}{128} = 0.0078125.
\end{equation}
The three gauge couplings are then expressed as integer multiples:
\begin{align}
\alpha_{\rm EM} &= 1 \times \Phi_T^2 = \frac{1}{128}, \\
\alpha_W        &= 4 \times \Phi_T^2 = \frac{4}{128}, \\
\alpha_S        &= 15 \times \Phi_T^2 = \frac{15}{128}.
\end{align}
The integers $\{1, 4, 15\}$ are base coefficients. They are assigned,
not rigorously derived from $SU(3) \times SU(2)$ group theory;
they are geometric coefficients on the assigned unified base $\Phi_T^2$.

The coupling ratios are therefore:
\begin{equation}
\alpha_S : \alpha_W : \alpha_{\rm EM} = 15 : 4 : 1.
\end{equation}

\section{The $\alpha_S/\alpha_W = 3.75$ Prediction}

\subsection{Derivation}

From the integer coefficients:
\begin{equation}
\frac{\alpha_S}{\alpha_W} = \frac{15}{4} = 3.75.
\end{equation}
This ratio is independent of the absolute value of $\Phi_T^2$, making it
a prediction of the framework. However, it depends on the assigned integers
$\{4, 15\}$ rather than being derived from first principles.

\subsection{Experimental Comparison}

\begin{table}[h]
\caption{Comparison of gauge coupling ratios.}
\begin{ruledtabular}
\begin{tabular}{lccc}
Quantity & Framework & Experiment (PDG 2024) & Deviation \\ \hline
$\alpha_S/\alpha_W$ & 3.7500 & $3.49$ & 7.4\% \\
$\alpha_S/\alpha_{\rm EM}$ & 15.00 & 15.09 & 0.6\% \\
$\alpha_W/\alpha_{\rm EM}$ & 4.000 & 4.32 & 8.1\% \\
\end{tabular}
\end{ruledtabular}
\end{table}

The predicted ratio $\alpha_S/\alpha_W = 3.75$ differs from experiment
by 7.4\%. The value $3.7959$ used a circular definition $\alpha_W=4\alpha_{\rm EM}$. The framework
does not predict the absolute values of $\alpha$ (involving $\Phi_T^2 = 1/128$
vs.\ the experimental $1/137$), so $\alpha_S/\alpha_{\rm EM}$ matches (0.6\%) but $\alpha_W/\alpha_{\rm EM}$ deviates by 8.1\%.

\subsection{Testability}

The prediction $\alpha_S/\alpha_W = 3.75$ is falsifiable:
\begin{itemize}
\item If future experiments at LHC, ILC, or CEPC measure
  $\alpha_S/\alpha_W = 3.75$ constant across scales, the framework is supported.
\item If the measurement diverges significantly from 3.75, the framework
  is falsified.
\end{itemize}
The Standard Model does not predict a specific integer value for this ratio;
our framework makes a falsifiable prediction, though based on assigned
coefficients rather than first-principles derivation.

\section{Honest Limitations}

\subsection{H1: The Fine Structure Constant}

The framework assigns $\Phi_T^2 = 1/128$, giving $\alpha_{\rm geom} = 1/128 =
0.0078125$. The experimental value is $\alpha_{\rm exp} = 1/137.036 = 0.00729735$,
a deviation of 7.06\%.

One-loop running corrections were investigated; they increase $\alpha$
(running toward $1/113.9$), but the experimental $\alpha(M_Z) = 1/128.9$
is smaller than $1/128$. The running makes the discrepancy worse, not better.

\textbf{The framework can accommodate but cannot predict $\alpha$ from first
principles.} This is an honest, unsolved problem.

\textbf{Note on previous claims:} Earlier versions of this work claimed
$\kappa^{28}\cdot S_{\rm dS}/S_{\rm BH} = 10^{-98}$ as an exact coincidence.
This claim was erroneous; the correct product is approximately $10^{-53}$
using standard formulas for de Sitter entropy $S_{\rm dS} = A/(4\ell_P^2)$.
This erroneous claim is hereby explicitly retracted.

\subsection{H2: Gravity as a Cosmological Quantity}

The gravitational coupling $G$ depends on the cosmological boundary condition:
$G$ is cosmological, not purely topological. The framework achieves
``3+1'' unification (three forces unified geometrically, gravity separate)
rather than full ``4'' unification.

\subsection{H3: Particle Masses}

The framework cannot calculate the masses of electrons, quarks, or neutrinos.
The Higgs mechanism and Yukawa couplings are not yet derived. This is a major
limitation, honestly acknowledged.

\subsection{H4: Input Parameters}

The parameters $\kappa$, $\tau$, and $\Phi_T$ are inputs, not derived from a deeper
principle. They are treated as ``constants of nature'' within the framework.
A derivation from 32-dimensional geometry or string theory remains as future work.

\section{Discussion}

\subsection{What the Framework Achieves}

\begin{itemize}
\item A geometric picture: four forces as perpendicular modes of light.
\item A novel interpretation of why four forces exist (and not three or five).
\item An explanation for the geometric meaning of $\alpha = \tau/\kappa$.
\item A weak falsifiable prediction: $\alpha_S/\alpha_W = 3.75$ (7.4\% deviation from independent experiment; best match is $\alpha_S/\alpha_{\rm EM}=15$ at 0.6\%).
\end{itemize}

\subsection{Relationship to Existing Theories}

This framework is not in competition with string theory or loop quantum
gravity. It provides a complementary geometric picture at the level of
gauge coupling unification. The value is primarily philosophical and
interpretational.

\subsection{What Would Make This a Theory}

To transition from ``geometric interpretation'' to ``theory'':
\begin{enumerate}
\item Derive $\kappa$ and $\tau$ from first principles.
\item Derive the unified seed $\Phi_T^2$ from geometry.
\item Calculate particle masses.
\item Confirm $\alpha_S/\alpha_W = 3.75$ at 0.1\% precision.
\end{enumerate}
Without these, the framework remains a geometric interpretation.

\section{Conclusion}

We have proposed that the four fundamental forces correspond to four
perpendicular oscillation modes of a light-speed helix structure in space.
The framework derives the gauge coupling ratio $15 : 4 : 1$ from assigned integer
coefficients on a unified base, yielding the falsifiable prediction
$\alpha_S/\alpha_W = 3.75$ (independent experiment: $3.49$, deviation 7.4\%; the value $3.7959$ was circular).

The framework is approximately 65\% complete as a geometric interpretation.
Key unsolved problems are honestly acknowledged: the absolute value of $\alpha$
(7.06\% deviation, H1), the non-calculation of particle masses (H3), and the
input status of $\kappa$, $\tau$, and $\Phi_T$ (H4). Previous claims of exact
numerical coincidences ($\kappa^{28}\cdot S_{\rm dS}/S_{\rm BH} = 10^{-98}$)
were erroneous and are hereby corrected.

We invite collaboration with phenomenologists and experimentalists to test
the $\alpha_S/\alpha_W = 3.75$ prediction at higher precision, and with
theorists working on gauge unification and quantum gravity to develop these
ideas further.

\begin{acknowledgments}
We acknowledge valuable discussions on gauge coupling unification and
geometric approaches to quantum gravity.
\end{acknowledgments}

\begin{thebibliography}{99}

% GUT history
 bibitem{GeorgiGlashow}
 H. Georgi and S. L. Glashow,
 ``Unity of All Elementary Particle Forces,''
 Phys. Rev. Lett.\ {\bf 32}, 438 (1974).

% PDG
 bibitem{PDG2024}
 R. L. Workman et al.\ (Particle Data Group),
 ``Review of Particle Physics,''
 Prog. Theor. Exp. Phys.\ {\bf 2024}, 083C01 (2024).

% Standard QFT text
 bibitem{Peskin}
 M. E. Peskin and D. V. Schroeder,
 ``An Introduction to Quantum Field Theory,''
 Addison-Wesley (1995).

% Geometric unification
 bibitem{Kaluza}
 T. Kaluza,
 ``On the Unification Problem in Physics,''
 Sitzungsber. Preuss. Akad. Wiss. Berlin (Math. Phys.)\ {\bf 1921}, 966 (1921).

% String theory
 bibitem{GSW}
 M. B. Green, J. H. Schwarz, and E. Witten,
 ``Superstring Theory,''
 Cambridge University Press (1987).

% LQG
 bibitem{Rovelli}
 C. Rovelli,
 ``Quantum Gravity,''
 Cambridge University Press (2004).

% SUSY GUT couplings
 bibitem{SUSYGUT}
 S. Raby,
 ``Grand Unified Theories,''
 in \textit{Particle Physics: A Comprehensive Description},
 Springer (2021).

\end{thebibliography}

\end{document}
